\uuid{vzgP}
\titre{Notion d'équation différentielle}
\auteur{bodin ; croizat ; sacre}
\organisation{exo7}
\licence{CC BY-NC-SA 4.0}
\datecreate{2022-03-08}
\niveau{L1}
\module{Analyse}
\chapitre{Equation différentielle}
\sousChapitre{Résolution d'équation différentielle du premier ordre}
\difficulte{2}
\theme{Notion d'équation différentielle}
\exo7qcmid{hilisit.equations-differentielles.019}

\contenu{
\begin{question}
\qclassification{225.01}
Soit $f$ une solution de l'équation différentielle $y'=y^2 + 1$.
Quelles sont les affirmations vraies sur la fonction $f$ ?
\begin{answers}
  \good{$f$ est une fonction croissante.}
  \bad{$f$ est une fonction décroissante.}
  \good{$f'$ est une fonction positive.}
  \bad{$f$ peut être une fonction constante.}
\end{answers}
\begin{explanations}
Si $f$ est solution de l'équation $y'=y^2 + 1$, alors on a $f'(x) = f^2(x) + 1$ et donc $f'(x) \geq 1 > 0$. Ainsi $f'$ est strictement positive, et par conséquent $f$ est strictement croissante.
\end{explanations}
\end{question}
}
